%
% --- inline annotations
%
% \usepackage[dvipsnames]{xcolor}
% \newcommand{\red}[1]{{\color{red}#1}}
% \newcommand{\todo}[1]{{\color{red}#1}}
% \newcommand{\TODO}[1]{\textbf{\color{red}[TODO: #1]}}
% --- disable by uncommenting  
% \renewcommand{\TODO}[1]{}
% \renewcommand{\todo}[1]{#1}

% \usepackage{xcolor}

\definecolor{lightindigo}{RGB}{242, 244, 255}
\definecolor{lightanthracite}{RGB}{252, 252, 255}
\definecolor{blueanthracite}{RGB}{61,71,94}
\definecolor{dark}{RGB}{31,31,31}
\definecolor{indigo}{RGB}{63, 81, 181}
\definecolor{red}{RGB}{210, 40, 95} 
\definecolor{pink}{RGB}{236, 64, 122}
\definecolor{green}{RGB}{46, 182, 125}
\definecolor{blue}{RGB}{66, 133, 244}
\definecolor{yellow}{RGB}{236, 178, 46}
\definecolor{anthracite}{RGB}{10, 4, 6}
\definecolor{gold}{RGB}{182, 131, 76}
\definecolor{lightgrey}{RGB}{128, 128, 128}
\definecolor{primary}{RGB}{40, 53, 147}
\definecolor{confcolor}{RGB}{63,81,181}
\definecolor{lightpink}{RGB}{252, 228, 236}
\definecolor{lightyellow}{RGB}{255, 253, 231}
\definecolor{teal}{RGB}{56,173,169}
\definecolor{beige}{RGB}{251,241,216}
\definecolor{temp}{RGB}{16,11,38}
\definecolor{teal}{RGB}{10, 126, 140}


\usepackage{graphicx}
% \usepackage{subfigure}
\usepackage{subcaption}
\usepackage{booktabs} % for professional tables
\usepackage{microtype}
% \usepackage{graphicx}
% \usepackage{subfigure}
% \usepackage{booktabs} % for professional tables


\usepackage{multirow}
\usepackage{makecell}
\usepackage{wrapfig} % For wrapping figures and tables
\usepackage{enumitem}

\usepackage{pifont}
\newcommand{\cmark}{\textcolor{green}{\ding{51}}} % Check mark symbol
\newcommand{\xmark}{\textcolor{red}{\ding{55}}} % Cross mark symbol
% \usepackage{stfloats}
\newcommand{\capprox}{\textcolor{orange}{\approx}}


% For theorems and such
\usepackage{amsmath}
\usepackage{amssymb}
\usepackage{mathtools}
\usepackage{amsthm}
\usepackage[normalem]{ulem}
% \usepackage{algorithm}
% \usepackage{algpseudocode}
\usepackage{algorithm}
% \usepackage{algorithmicx}
\usepackage{algpseudocode}

\usepackage{hyperref}
%\usepackage[capitalize,noabbrev]{cleveref}
\usepackage[capitalize]{cleveref}

%%%%%%%%%%%%%%%%%%%%%%%%%%%%%%%%
% THEOREMS
%%%%%%%%%%%%%%%%%%%%%%%%%%%%%%%%
\theoremstyle{plain}
\newtheorem{theorem}{Theorem}[section]
\newtheorem{proposition}[theorem]{Proposition}
\newtheorem{lemma}[theorem]{Lemma}
\newtheorem{corollary}[theorem]{Corollary}
\newtheorem{definition}[theorem]{Definition}
\newtheorem{assumption}[theorem]{Assumption}
\theoremstyle{definition}
\theoremstyle{remark}
\newtheorem{remark}[theorem]{Remark}


% Todonotes is useful during development; simply uncomment the next line
%    and comment out the line below the next line to turn off comments
\usepackage[disable,textsize=tiny]{todonotes}
% \usepackage[textsize=tiny]{todonotes}

\newif\ifshowcomments
% \showcommentstrue % Comment this line to disable comments

% Define the comment commands
\ifshowcomments
\newcommand{\thib}[1]{{\color{blue}[thib]: #1}}
\newcommand{\agus}[1]{{\color{blue}[agus]: #1}}
\newcommand{\fel}[1]{{\color{pink}[fel]: #1}}
\newcommand{\franck}[1]{{\color{red}[franck]: #1}}
\else
\newcommand{\thib}[1]{}
\newcommand{\agus}[1]{}
\newcommand{\fel}[1]{}
\newcommand{\franck}[1]{}
\fi

% \newcommand{\libname}{at \href{https://github.com/thib-s/orthogonium}{ \texttt{https://github.com/thib-s/orthogonium}}}
\newcommand{\libname}{\href{https://github.com/deel-ai/orthogonium}{here}}
\newcommand{\fulllibname}{\href{https://github.com/deel-ai/orthogonium}{Orthogonium}}
% \newcommand{\libname}{ \href{https://anonymous.4open.science/r/orthogonium-DE07}{here}}

\newcommand{\RR}{\mathbb{R}}
\newcommand{\noyau}{\mathbf{K}}
\newcommand{\Kbf}{\noyau}
\newcommand{\KbfS}{\RR^{c_o\times c_i\times k\times k}}   
\newcommand{\KbbfS}{\RR^{c_o\times c_i\times k_1\times k_2}}  
% version 2D
% \newcommand{\KbfS}{\RR^{c_o\times \frac{c_i}{g}\times k_1\times k_2}}    % version 2D
\newcommand{\cKbf}{\text{conv}_{\noyau}}

\newcommand{\Kk}{\mathcal{K}}
%\newcommand{\KkS}{\RR^{c_o\frac{h}{s}\frac{w}{s} \times c_ihw}} 
\newcommand{\KkS}{\RR^{c_ohw \times c_ihw}} 

\newcommand{\flashlipschitz}{turbo-BCOP}

\DeclareMathOperator{\Vect}{Vect}

\newcommand{\Bconv}{Block-convolution}
\newcommand{\compat}[2]{#1\Join#2}

% \newcommand{\mathbconv}{\divideontimes} % or \botdot or \circledcirc \circledast \divideontimes
\newcommand{\mathbconv}{\circledast} % or \botdot or \circledcirc \circledast \divideontimes
\newcommand{\kernel}[1]{\mathbf{#1}}
\newcommand{\kernelid}{\kernel{I}}
\newcommand{\toep}[1]{\mathcal{#1}}
\newcommand{\toepvec}[1]{\bar{#1}}
\newcommand{\toepstride}[1]{\toep{S}_{#1}}
\newcommand{\toepconv}[2]{\toep{S}_{#2}\toep{#1}}
\newcommand{\toepid}{\toep{I}}
\newcommand{\convcode}[3]{\mathtt{conv}_{#1}(#2, \mathtt{stride}=#3)}
\newcommand{\shortconvcode}[3]{\mathtt{conv}_{#1}(#2, #3)}
\newcommand{\convtrcode}[3]{\mathtt{ConvTranspose}_{#1}(#2, \mathtt{stride}=#3)}
\newcommand{\convker}[1]{\star_{#1}}
\newcommand{\RRkernel}[4]{\RR^{#1 \times #2 \times #3 \times #4}}

% \widowpenalty=50     % default is 150, lower it to allow widows
% \clubpenalty=50      % default is 150, lower it to allow orphans
% \brokenpenalty=50    % default is 1000, allows breaks at hyphenated lines
% \displaywidowpenalty=50  % default is 50, allows widows in displayed equations
% \predisplaypenalty=50    % default is 10000, allows page breaks before displayed equations
% % \setlength{\parskip}{-1pt}      % remove extra space between paragraphs
% % \setlength{\parindent}{1em}    % reduce paragraph indentation if necessary
